\chapter{Lezione 6}

\textbf{Argomenti:} Risposta del modello HH a correnti costanti, treni di spike, curva frequenza-corrente (F-I), neuroni di Tipo I e Tipo II, adattamento della frequenza di scarica (SFA), canali HVA calcio-dipendenti, dinamica intracellulare del Ca$^{2+}$, canali K$^+$ Ca$^{2+}$-dipendenti, trasmissione sinaptica chimica, esocitosi, EPSP/IPSP, stochasticità sinaptica, analisi quantale, modello cinetico della conduttanza sinaptica, filtro sinaptico passa-basso

\section{Risposta del Modello HH a Correnti Costanti}

Il professore riprende dal modello di \textbf{Hodgkin-Huxley (HH)}, che descrive l'eccitabilità della membrana neuronale attraverso le variabili di gate $m$, $h$ ed $n$ per i canali Na$^+$ voltaggio-dipendenti (attivazione e inattivazione) e K$^+$ voltaggio-dipendenti (attivazione ritardata). L'obiettivo di questa sezione è comprendere come il neurone risponde a una \textbf{corrente costante a gradino} (step current), che modella in modo semplificato il \textbf{bombardamento sinaptico} continuo che il neurone riceve in vivo dai neuroni della rete.

La corrente iniettata $I_e$ (corrente esterna) simula la somma di tutti gli input sinaptici eccitatori e inibitori che arrivano sul dendrite e sul soma. In laboratorio, questa procedura è realizzata sperimentalmente con l'\textbf{iniezione di corrente intracellulare} attraverso un elettrodo.

\subsubsection{L'Equazione Sub-soglia del Potenziale di Membrana}

Quando la corrente iniettata $I_e$ è sufficientemente piccola, il neurone non supera la soglia di produzione di potenziali d'azione. In questo \textbf{regime lineare sub-soglia}, i canali Na$^+$ e K$^+$ voltaggio-dipendenti non si aprono significativamente: la dinamica del potenziale di membrana $V$ è dominata dal solo canale di perdita (leak) con conduttanza $G_L$ e potenziale di inversione $E_L \approx E_m = -65$ mV (nel modello HH standard).

L'equazione che governa questo regime è:

\begin{equation}
C_m \frac{dV}{dt} = -G_L(V - E_L) + \frac{I_e}{A}
\end{equation}

dove:
\begin{itemize}
    \item $C_m$ è la capacità di membrana per unità di area $[\mu\text{F/cm}^2]$,
    \item $G_L$ è la conduttanza di perdita per unità di area $[\text{mS/cm}^2]$,
    \item $E_L$ è il potenziale di inversione del leak $[\text{mV}]$,
    \item $\frac{I_e}{A}$ è la densità di corrente iniettata $[\mu\text{A/cm}^2]$ (corrente totale divisa per l'area della membrana).
\end{itemize}

\subsubsection{Deviazione Stazionaria dal Potenziale di Riposo}

In condizioni stazionarie ($dV/dt = 0$), il potenziale di membrana raggiunge un nuovo valore di equilibrio $V_{ss}$ spostato rispetto al valore di riposo $E_m = -65$ mV. La \textbf{deviazione stazionaria} $\Delta V$ vale:

\begin{equation}
\Delta V = V_{ss} - E_m = \frac{I_e / A}{G_L}
\end{equation}

Questa è semplicemente la legge di Ohm applicata al circuito RC della membrana. Il neurone si comporta come un filtro passa-basso del primo ordine con costante di tempo $\tau_m = C_m / G_L$ (la \textbf{costante di tempo di membrana}).

\subsubsection{Oscillazioni Smorzate nell'Avvicinamento all'Equilibrio}

Un fenomeno caratteristico e non banale del modello HH è che, durante la transizione verso il nuovo potenziale stazionario, il $V_m$ \textbf{non si avvicina monotonicamente} all'equilibrio, bensì presenta \textbf{oscillazioni smorzate} (damped oscillations). Il sistema HH è, vicino al punto di equilibrio, un oscillatore sottoammortizzato: il potenziale di membrana oscilla attorno a $V_{ss}$ con ampiezza decrescente esponenzialmente prima di stabilizzarsi.

\begin{tcolorbox}[colback=gray!5, colframe=gray!40, arc=2mm, boxrule=0.5pt]
Queste oscillazioni smorzate sono un primo indizio della \textbf{struttura di biforcazione} del sistema HH: il punto fisso è uno spirale stabile (stable focus) e il sistema si trova vicino a una biforcazione di Hopf. Questo è connesso alla classificazione dei neuroni in Tipo I e Tipo II che vedremo tra poco.
\end{tcolorbox}



\section{Treni di Spike}

\subsubsection{Soglia e Produzione di un Singolo Spike}

Quando la corrente iniettata viene raddoppiata a $I_e \approx 4 \, \mu\text{A/cm}^2$, il sistema supera la \textbf{soglia di eccitabilità}: si produce un singolo \textbf{potenziale d'azione} (spike). Il meccanismo è il seguente:

\begin{enumerate}
    \item La corrente capacitiva iniziale carica rapidamente la membrana (fase di ramping di $V_m$).
    \item Quando $V_m$ supera la soglia $V_{th} \approx -58.5$ mV (nel modello HH), i canali Na$^+$ voltaggio-dipendenti si aprono con retroazione positiva (overshoot).
    \item Il potenziale raggiunge il \textbf{potenziale di picco} (peak) vicino al potenziale di inversione del sodio: $E_{Na} \approx +50$ mV.
    \item Segue la \textbf{fase di ripolarizzazione} (apertura dei canali K$^+$ ritardati e inattivazione dei canali Na$^+$) e l'\textbf{iperpolarizzazione} (la cosiddetta \textbf{undershoot} o \textbf{periodo refrattario}).
    \item Il potenziale ritorna lentamente al valore di riposo $E_m$.
\end{enumerate}


\subsubsection{Dal Singolo Spike al Treno di Spike}

A corrente $I_e = 8 \, \mu\text{A/cm}^2$: si produce un \textbf{treno di spike} (spike train). La corrente iniettata è ora sufficiente a mantenere il potenziale di membrana in un regime che supera \textbf{ripetutamente e periodicamente} la soglia, generando una sequenza di potenziali d'azione.

Il meccanismo è il seguente: dopo ogni spike, durante il periodo refrattario il $V_m$ scende sotto soglia (iperpolarizzazione); poi la corrente iniettata fa ricrescere lentamente il $V_m$ (ramping) fino a raggiungere nuovamente la soglia, producendo un nuovo spike. Il processo si ripete indefinitamente a corrente costante.

A $I_e = 16 \, \mu\text{A/cm}^2$: il treno è ancora più denso, perché $dV/dt \propto I_e$ è più elevato $\rightarrow$ il ramping è più veloce $\rightarrow$ la soglia viene raggiunta in minor tempo $\rightarrow$ periodo tra spike più breve.


\subsection{Frequenza di Scarica e Inter-Spike Interval (ISI)}

Durante un treno di spike periodico, il parametro più importante è l'\textbf{Inter-Spike Interval (ISI)}, cioè il tempo tra due spike consecutivi:

\begin{equation}
T_{ISI} = t_{f+1} - t_f
\end{equation}

dove $t_f$ è il tempo del $f$-esimo spike e $t_{f+1}$ è il tempo del successivo.

La \textbf{frequenza di scarica} (firing rate) $\nu$ è l'inverso dell'ISI:

\begin{equation}
\nu = \frac{1}{T_{ISI}}
\end{equation}

Si misura in Hz (scariche al secondo). Per un treno perfettamente periodico, $\nu$ è costante; in presenza di adattamento o rumore, $\nu$ varia nel tempo.



L'ISI dipende direttamente dall'intensità della corrente iniettata:
\begin{itemize}
    \item \textbf{Alta corrente} $\Rightarrow$ ramping più rapido $\Rightarrow$ $T_{ISI}$ più corto $\Rightarrow$ $\nu$ più alta.
    \item \textbf{Bassa corrente (ma sopra soglia)} $\Rightarrow$ ramping lento $\Rightarrow$ $T_{ISI}$ lungo $\Rightarrow$ $\nu$ bassa.
\end{itemize}

La relazione tra $I_e$ e $\nu$ definisce la \textbf{curva F-I} o curva guadagno del neurone.


\subsection{La Curva F-I del Modello HH}

La \textbf{curva frequenza-corrente} (F-I curve, anche chiamata curva guadagno o curva di trasferimento) mappa la corrente iniettata $I_e$ nella frequenza di scarica stazionaria $\nu$:

\begin{equation}
\nu = f(I_e)
\end{equation}

Per il modello HH, la curva F-I ha le seguenti proprietà:
\begin{itemize}
    \item Per $I_e$ piccola (sotto soglia): $\nu = 0$ (nessun firing).
    \item Esiste una \textbf{corrente di reobase} $I_{rh}$ (rheobase current): la corrente minima necessaria per produrre spike. Sotto $I_{rh}$, il sistema HH non oscilla.
    \item A $I_e = I_{rh}$, il firing rate \textbf{salta discontinuamente} a un valore $\nu_{min} > 0$: non si può avere firing a frequenza arbitrariamente bassa.
\end{itemize}

Questo comportamento è la firma dei \textbf{neuroni di Tipo II}.

\subsubsection{Neuroni di Tipo II (Discontinuous F-I Curve)}

Un \textbf{neurone di Tipo II} è caratterizzato da una curva F-I con salto discontinuo alla reobase. Il modello HH è l'esempio prototipico di neurone di Tipo II.

La caratteristica matematica sottostante è che il punto fisso del sistema HH perde stabilità attraverso una \textbf{biforcazione di Hopf subcritica}, nella quale il ciclo limite (oscillazione periodica) appare a frequenza finita. Questo spiega il salto: il neurone passa direttamente dalla quiescenza a una frequenza minima $\nu_{min}$.

\begin{tcolorbox}[colback=gray!5, colframe=gray!40, arc=2mm, boxrule=0.5pt]
La biforcazione di Hopf subcritica implica che il sistema HH è \textbf{bistabile} in un piccolo intervallo di corrente attorno alla reobase: coesistono il punto fisso stabile e il ciclo limite stabile. Questo ha conseguenze importanti per la risposta del neurone a stimoli transitori.
\end{tcolorbox}

\subsubsection{Neuroni di Tipo I (Continuous F-I Curve)}

I \textbf{neuroni di Tipo I} hanno invece una curva F-I \textbf{continua}: la frequenza di firing parte da zero alla reobase e cresce in modo continuo (anche se la derivata $d\nu/dI_e$ può essere discontinua o molto elevata alla soglia). Un neurone di Tipo I può perciò scaricare a frequenze arbitrariamente basse, basta avere una corrente di poco superiore alla reobase.

Il meccanismo sottostante è una \textbf{biforcazione SNIC (Saddle-Node on Invariant Circle)}: il punto fisso collide con un punto di sella su un ciclo invariante. Il ciclo limite nasce con frequenza zero e cresce continuamente.

Esempio biologico: i neuroni della corteccia di tipo RS (Regular Spiking) si comportano approssimativamente come neuroni di Tipo I.

\subsubsection{Classificazione come Prima Tassonomia dei Neuroni}

Questa distinzione Tipo I / Tipo II rappresenta la \textbf{prima classificazione dei neuroni} a seconda del loro \textbf{pattern di eccitabilità}. Fu introdotta da \textbf{Rinzel e Ermentrout} negli anni '80-'90, sulla base dell'analisi delle biforcazioni dei sistemi dinamici neurali. Ha implicazioni importanti per la codifica dell'informazione: un neurone di Tipo I può codificare segnali deboli attraverso la modulazione fine della frequenza; un neurone di Tipo II risponde in modo più tutto-o-nulla.


\section{Dinamica del Calcio}


Negli anni '30, \textbf{Sir Edgar Douglas Adrian} osservò sperimentalmente che i neuroni sensoriali non scaricano a frequenza costante in risposta a uno stimolo costante: la frequenza di scarica \textbf{decresce progressivamente nel tempo}, anche a corrente iniettata costante. Questo fenomeno è chiamato \textbf{adattamento della frequenza di scarica} (Spike Frequency Adaptation, SFA).

Adrian ricevette il Premio Nobel per la Fisiologia o Medicina nel 1932 per le sue scoperte sulla funzione dei neuroni, tra cui questo fenomeno.

Il modello HH standard (con soli canali Na$^+$ e K$^+$ voltaggio-dipendenti) \textbf{non predice} l'SFA: a corrente costante, il modello HH produce un treno di spike a frequenza perfettamente costante. Per spiegare l'SFA è necessario introdurre canali ionici aggiuntivi.

\subsubsection{Canali HVA (High-Voltage Activated) del Calcio}

La chiave dell'SFA è l'esistenza di \textbf{canali ionici Ca$^{2+}$ voltaggio-dipendenti} di tipo \textbf{HVA} (High-Voltage Activated). Questi canali:

\begin{itemize}
    \item Si aprono quando il potenziale di membrana è \textbf{molto elevato}, tipicamente durante la fase ascendente di un potenziale d'azione ($V_m > -20$ mV circa).
    \item Rimangono sostanzialmente chiusi al potenziale di riposo.
    \item Permettono un rapido \textbf{influsso di Ca$^{2+}$} nella cellula durante ogni spike.
\end{itemize}

La concentrazione di calcio extracellulare $[\text{Ca}^{2+}]_o \approx 2$ mM è molto maggiore di quella intracellulare (a riposo: $[\text{Ca}^{2+}]_i \approx 100$ nM $= 10^{-7}$ M), quindi il gradiente elettrochimico favorisce fortemente l'ingresso del calcio.

\begin{tcolorbox}[colback=gray!5, colframe=gray!40, arc=2mm, boxrule=0.5pt]
 I canali HVA del calcio esistono in diversi sottotipi (L, N, P/Q, R), con caratteristiche di cinetica e localizzazione subcellulare diverse. Ai fini del modello di SFA che trattiamo, interessa solo che si aprano durante gli spike e permettano l'ingresso di Ca$^{2+}$.
\end{tcolorbox}


\subsubsection{Equazione Differenziale per la Concentrazione di Calcio}

La concentrazione di calcio intracellulare $[\text{Ca}^{2+}]_i$ (abbreviata $[\text{Ca}]_i$ nelle equazioni) varia dinamicamente durante il treno di spike. La sua evoluzione temporale è descritta da un'equazione del primo ordine:

\begin{equation}
\frac{d[\text{Ca}]_i}{dt} = -\frac{[\text{Ca}]_i}{\tau_{Ca}} + \alpha_{Ca} \sum_{f} \delta(t - t_f)
\end{equation}


\begin{itemize}
    \item $\tau_{Ca} \approx 150$ ms è la \textbf{costante di tempo di rilassamento del Ca$^{2+}$} verso il valore di equilibrio (bassissimo);
    \item $\alpha_{Ca}$ è l'\textbf{ampiezza del salto} di $[\text{Ca}]_i$ per ogni spike;
    \item $t_f$ sono i tempi dei potenziali d'azione;
    \item $\delta(t - t_f)$ è la delta di Dirac, che modella l'apertura istantanea dei canali HVA durante lo spike.
\end{itemize}

Il termine $-[\text{Ca}]_i / \tau_{Ca}$ descrive il rilassamento esponenziale verso la concentrazione basale (effettivamente $\approx 0$ nel modello semplificato), dovuto ai \textbf{trasportatori del calcio} (pompa Ca$^{2+}$-ATPasi, scambiatore Na$^+$/Ca$^{2+}$, tamponamento intracellulare).

\subsubsection{Accumulo di Ca$^{2+}$ durante il Treno di Spike}

Durante un treno di spike con periodo $T_{ISI}$, ogni spike aggiunge un salto $\alpha_{Ca}$ a $[\text{Ca}]_i$. Se il decadimento nell'ISI non è completo (cioè se $T_{ISI} \ll \tau_{Ca}$), la concentrazione di calcio si \textbf{accumula progressivamente}:

\begin{itemize}
    \item Dopo il primo spike: $[\text{Ca}]_i \approx \alpha_{Ca}$.
    \item Dopo il secondo spike: $[\text{Ca}]_i \approx 2\alpha_{Ca} - \alpha_{Ca}(1 - e^{-T_{ISI}/\tau_{Ca}}) = \alpha_{Ca}(1 + e^{-T_{ISI}/\tau_{Ca}})$.
    \item Asintoticamente, $[\text{Ca}]_i$ si stabilizza su un valore medio $\langle[\text{Ca}]_i\rangle = \alpha_{Ca} \tau_{Ca} / T_{ISI} \cdot (1 - e^{-T_{ISI}/\tau_{Ca}})^{-1}$.
\end{itemize}

Poiché $\tau_{Ca} \approx 150$ ms è molto maggiore di un tipico ISI (ad es. 20–50 ms), il Ca$^{2+}$ si accumula efficacemente spike dopo spike.

\subsubsection{Canali K$^+$ Ca$^{2+}$-Dipendenti e la Variabile di Gate $n_\infty$}

Il calcio accumulato attiva una \textbf{corrente di K$^+$ aggiuntiva} attraverso \textbf{canali K$^+$ Ca$^{2+}$-dipendenti} (SKCa: Small-conductance Calcium-activated Potassium channels). Questi canali sono \textbf{diversi} dai canali K$^+$ voltaggio-dipendenti del modello HH standard: si aprono in risposta alla concentrazione intracellulare di Ca$^{2+}$, non al potenziale di membrana.

La variabile di gate $n$ per questi canali ha un'asintoto sigmoidale:

\begin{equation}
n_\infty([\text{Ca}]) = \frac{[\text{Ca}]}{[\text{Ca}] + K_2}
\end{equation}

dove $K_2$ è una costante di dissociazione:
\begin{itemize}
    \item Quando $[\text{Ca}]_i \approx 0$ (a riposo): $n_\infty \approx 0$ $\rightarrow$ canali chiusi.
    \item Quando $[\text{Ca}]_i \gg K_2$: $n_\infty \to 1$ $\rightarrow$ canali completamente aperti.
\end{itemize}

La costante di tempo di questa variabile di gate è $\tau_n \approx 10$ ms, che è molto \textbf{più rapida} di $\tau_{Ca} \approx 150$ ms. Questo significa che $n$ insegue quasi istantaneamente $[\text{Ca}]_i$: si può applicare l'\textbf{approssimazione adiabatica} $n \approx n_\infty([\text{Ca}]_i)$.

\subsubsection{Corrente K$^+$ Ca$^{2+}$-Dipendente}

La corrente $K^+$ Ca$^{2+}$-dipendente è:

\begin{equation}
I_{K,Ca} = G_{K,Ca} \cdot n_\infty([\text{Ca}]_i) \cdot (V - E_K)
\end{equation}

dove $E_K \approx -76$ mV (potenziale di inversione del K$^+$). Poiché durante il normale firing $V > E_K$, questa corrente è \textbf{uscente} (positiva per convenzione elettrofisiologica) e quindi \textbf{iperpolarizzante}: tende ad allontanare il potenziale di membrana dalla soglia.


\subsubsection{Come la Corrente K$^+$ Ca$^{2+}$-Dipendente Allunga l'ISI}

Il meccanismo dell'\textbf{adattamento della frequenza di scarica} è ora chiaro: la corrente $I_{K,Ca}$ (negativa = iperpolarizzante) si \textbf{somma algebricamente} alla corrente iniettata $I_e$, riducendo la corrente netta disponibile per caricare la membrana:

\begin{equation}
I_{netto}(t) = \frac{I_e}{A} - G_{K,Ca} \cdot n_\infty([\text{Ca}]_i(t)) \cdot (V - E_K)
\end{equation}

Poiché $[\text{Ca}]_i$ aumenta nel corso del treno, il termine $I_{K,Ca}$ cresce progressivamente nel tempo. Questo riduce $I_{netto}$, che a sua volta riduce la velocità di ramping del potenziale di membrana tra due spike ($dV/dt = I_{netto}/C_m$).

La conseguenza diretta è che \textbf{l'ISI si allunga progressivamente}: il potenziale ci mette più tempo a raggiungere la soglia, quindi la frequenza di scarica $\nu = 1/T_{ISI}$ \textbf{diminuisce nel tempo}.

\subsubsection{Andamento Temporale dell'Adattamento}

L'evoluzione dell'SFA ha due fasi temporali caratteristiche:

\begin{enumerate}
    \item \textbf{Fase rapida} (scala di $\tau_n \approx 10$ ms): ogni spike produce immediatamente un piccolo incremento di $I_{K,Ca}$, che allunga leggermente l'ISI successivo.
    \item \textbf{Fase lenta} (scala di $\tau_{Ca} \approx 150$ ms): il Ca$^{2+}$ si accumula lentamente, aumentando gradualmente il livello basale di $I_{K,Ca}$, con un ulteriore allungamento progressivo degli ISI.
\end{enumerate}

L'SFA porta il sistema verso un nuovo stato stazionario in cui l'ISI si è stabilizzato su un valore più lungo rispetto all'ISI iniziale. Se la corrente iniettata viene rimossa, $[\text{Ca}]_i$ decade con $\tau_{Ca} \approx 150$ ms e il sistema torna alla condizione di riposo.


\subsubsection{La Frequenza Massima Assoluta}

Esiste un limite fisico superiore alla frequenza di firing, determinato dal \textbf{periodo refrattario assoluto} $\tau_{ref}$. Durante il periodo refrattario, i canali Na$^+$ sono completamente inattivati (variabile $h \approx 0$) e il neurone è assolutamente incapace di produrre un nuovo spike, indipendentemente dall'entità della corrente iniettata.

La \textbf{frequenza massima di scarica} è quindi:

\begin{equation}
\nu_{max} = \frac{1}{\tau_{ref}} \approx \frac{1}{2 \times 10^{-3} \, \text{s}} = 500 \, \text{Hz}
\end{equation}

dove $\tau_{ref} \approx 2$ ms è il periodo refrattario assoluto tipico.

\subsubsection{Rilevanza Biologica}

In vivo, le frequenze di scarica tipiche nella corteccia cerebrale sono dell'ordine di \textbf{poche decine di Hz} (tipicamente 5–80 Hz), ben al di sotto di $\nu_{max} = 500$ Hz. Questo non è casuale: l'SFA è uno dei meccanismi che \textbf{impedisce} ai neuroni di avvicinarsi alla frequenza massima in condizioni fisiologiche normali.

\begin{tcolorbox}[colback=gray!5, colframe=gray!40, arc=2mm, boxrule=0.5pt]
 Se i neuroni corticali potessero scaricare a 500 Hz, il consumo energetico del cervello diventerebbe insostenibile (già a riposo il cervello consuma circa il 20\% dell'energia corporea totale). L'SFA agisce come un meccanismo di \textbf{risparmio energetico}: abbassa automaticamente il firing rate in risposta a stimoli prolungati.
\end{tcolorbox}

L'SFA ha anche un \textbf{significato computazionale}: favorisce la codifica dei \textbf{transitori} degli stimoli (onset) rispetto alla loro componente stazionaria. Un neurone con SFA risponde con un burst iniziale ad alta frequenza e poi adatta la sua scarica, "segnalando" l'inizio di un nuovo stimolo più che la sua durata.


\section{Trasmissione Sinaptica Chimica}

 Le correnti che il neurone riceve e che determinano il suo firing non sono generate direttamente dall'interno della cellula, ma arrivano da altri neuroni attraverso le \textbf{sinapsi}.

Nel sistema nervoso centrale dei mammiferi, la forma di comunicazione interneuronale più comune è la \textbf{trasmissione sinaptica chimica} (in opposizione alle \textbf{sinapsi elettriche} o gap junctions, che sono presenti in strutture più antiche filogeneticamente come il cervelletto e il talamo, ma anche in alcune regioni corticali).

\subsubsection{Morfologia della Sinapsi Chimica}

La struttura morfologica di una sinapsi chimica comprende tre componenti principali:

\begin{enumerate}
    \item \textbf{Bottone presinaptico} (presynaptic terminal o bouton): è il rigonfiamento terminale dell'assone del neurone presinaptico. Contiene migliaia di \textbf{vescicole sinaptiche} — piccole sfere membranose (diametro $\approx$ 40 nm) riempite di molecole di \textbf{neurotrasmettitore}.
    \item \textbf{Fessura sinaptica} (synaptic cleft): spazio extracellulare di circa \textbf{20 nm} che separa la membrana presinaptica da quella postsinaptica. È crucialmente stretto: i neurotrasmettitori possono diffondere attraverso di esso in pochi microsecondi.
    \item \textbf{Membrana postsinaptica}: si trova tipicamente sul \textbf{dendrite} (o sul soma) del neurone postsinaptico. Contiene alta densità di \textbf{recettori ionotropici} (canali ionici ligando-dipendenti) nella \textbf{densità postsinaptica} (PSD, PostSynaptic Density).
\end{enumerate}

\subsubsection{Neuroni Piramidali e Organizzazione Corticale}

I principali neuroni eccitatori della corteccia cerebrale sono i \textbf{neuroni piramidali}. Hanno un \textbf{dendrite apicale} orientato verticalmente verso la superficie corticale e numerosi \textbf{dendriti basali}. Le sinapsi sono distribuite sull'albero dendritico: le sinapsi apicali (provenienti da neuroni di altri strati corticali o aree remote) sono vicine alla superficie, mentre le sinapsi basali ricevono input locali.

\subsubsection{Il Ruolo del Calcio nel Processo Presinaptico}

Quando uno spike raggiunge il bottone presinaptico, la \textbf{depolarizzazione della membrana presinaptica} apre i canali Ca$^{2+}$ \textbf{HVA} presenti nel terminale. Il Ca$^{2+}$ extracellulare entra rapidamente nel bottone presinaptico. Questo evento è \textbf{indispensabile} per la trasmissione: in assenza di Ca$^{2+}$ extracellulare, gli spike presinaptici non producono rilascio di neurotrasmettitori.

Il Ca$^{2+}$ che entra nel bottone presinaptico si lega a proteine sensori del calcio (come la \textbf{sinaptotagmina}), inducendo la fusione delle vescicole sinaptiche con la membrana presinaptica (esocitosi).


\subsection{Meccanismo di Esocitosi}

\subsubsection{Le Proteine SNARE e il Meccanismo di Fusione}

Il meccanismo molecolare dell'esocitosi è mediato dalle \textbf{proteine SNARE} (Soluble NSF Attachment Protein Receptors):

\begin{itemize}
    \item \textbf{Synaptobrevin} (o VAMP): proteina SNARE sulla vescicola (v-SNARE).
    \item \textbf{Syntaxin} e \textbf{SNAP-25}: proteine SNARE sulla membrana presinaptica (t-SNARE).
\end{itemize}

Quando il Ca$^{2+}$ si lega alla \textbf{sinaptotagmina} (il sensore del calcio), si scatena un cambiamento conformazionale che permette alle proteine SNARE di formare un complesso a quattro eliche (SNARE complex). Questo complesso è un motore molecolare che "cuce" insieme la membrana vescicolare e quella presinaptica, permettendo la \textbf{fusione} e l'apertura di un poro di fusione (fusion pore).

\subsubsection{Rilascio dei Neurotrasmettitori nella Fessura}

Una volta aperto il poro di fusione, le molecole di \textbf{neurotrasmettitore} contenute nella vescicola (tipicamente migliaia di molecole per vescicola: ad es. $\sim$5000 molecole di glutammato per vescicola glutamatergica) vengono rilasciate nella fessura sinaptica per \textbf{diffusione}.

Date le dimensioni nanometriche della fessura (20 nm), la diffusione è estremamente rapida: le molecole di neurotrasmettitore raggiungono i recettori postsinaptici in pochi \textbf{microsecondi}.

\subsubsection{Apertura dei Canali Ionici Postsinaptici}

I neurotrasmettitori rilasciati si legano ai \textbf{recettori ionotropici} sulla membrana postsinaptica. I recettori ionotropici sono canali ionici \textbf{ligando-dipendenti}: si aprono quando una o più molecole di neurotrasmettitore si legano al loro sito di riconoscimento. Esempi principali:

\begin{table}[htbp]
    \centering
    \renewcommand{\arraystretch}{1.3}
    \begin{tabularx}{\textwidth}{|X|X|X|X|}
    \hline
    \textbf{Neurotrasmettitore} & \textbf{Recettore} & \textbf{Ione principale} & \textbf{Effetto} \\
    \hline
    Glutammato & AMPA & Na$^+$ & Depolarizzazione (EPSP) \\
    \hline
    Glutammato & NMDA & Na$^+$/Ca$^{2+}$ & Depolarizzazione (EPSP) \\
    \hline
    GABA & GABA$_A$ & Cl$^-$ & Iperpolarizzazione (IPSP) \\
    \hline
    Glicina & Recettore glicina & Cl$^-$ & Iperpolarizzazione (IPSP) \\
    \hline
    \end{tabularx}
\end{table}

Il lavoro sperimentale su questi meccanismi è stato premiato con il \textbf{Nobel per la Fisiologia o Medicina 1963}, assegnato a \textbf{John Eccles} (sinapsi), \textbf{Alan Hodgkin} e \textbf{Andrew Huxley} (potenziale d'azione).


\section{EPSP, IPSP e Chiusura Spontanea dei Canali}

\subsubsection{EPSP (Excitatory PostSynaptic Potential)}

Un \textbf{EPSP} (Potenziale Postsinaptico Eccitatorio) è la variazione di potenziale di membrana indotta dall'apertura di canali ionici eccitatori (tipicamente Na$^+$ o Na$^+$/K$^+$):

\begin{itemize}
    \item I canali AMPA si aprono $\rightarrow$ ingresso di Na$^+$ $\rightarrow$ depolarizzazione della membrana postsinaptica.
    \item Il potenziale di membrana si allontana dal potenziale di riposo verso valori meno negativi (più vicini a $E_{Na} \approx +50$ mV).
    \item L'ampiezza tipica di un singolo EPSP in un neurone corticale di mammifero è di $\sim 0.1$–$1$ mV (quindi molto al di sotto della soglia di firing $\sim 20$ mV sopra il riposo).
\end{itemize}

\subsubsection{IPSP (Inhibitory PostSynaptic Potential)}

Un \textbf{IPSP} (Potenziale Postsinaptico Inibitorio) è la variazione di potenziale opposta:

\begin{itemize}
    \item I canali GABA$_A$ si aprono $\rightarrow$ ingresso di Cl$^-$ (o uscita di K$^+$) $\rightarrow$ iperpolarizzazione della membrana postsinaptica.
    \item Il potenziale di membrana scende verso $E_{Cl}$ o $E_K$, che sono entrambi più negativi del potenziale di riposo.
    \item Questo rende il neurone postsinaptico più difficile da portare a soglia (inibizione).
\end{itemize}

\subsubsection{Decadimento dell'EPSP: Rising e Decaying Phase}

Una volta che i neurotrasmettitori iniziano a diffondere fuori dalla fessura sinaptica, la loro concentrazione nella fessura stessa diminuisce. Di conseguenza i recettori ionotropici postsinaptici \textbf{si chiudono spontaneamente} (per dissociazione del ligando), e la corrente postsinaptica decade.

L'andamento temporale di un EPSP presenta quindi due fasi:
\begin{enumerate}
    \item \textbf{Rising phase} (fase ascendente): durata $\sim 1$–$2$ ms, corrispondente all'apertura dei canali.
    \item \textbf{Decaying phase} (fase discendente): durata variabile (da $\sim 3$ ms per sinapsi veloci AMPA fino a centinaia di ms per sinapsi NMDA/metabotropiche), corrispondente alla chiusura spontanea dei canali.
\end{enumerate}




\subsubsection{L'Esperimento di Stimolazione Ripetuta}

Un esperimento fondamentale per comprendere la trasmissione sinaptica consiste nel stimolare ripetutamente il neurone presinaptico con spike identici (stessa ampiezza, stesso timing) e misurare la risposta postsinaptica. I risultati sono sorprendenti:

\begin{itemize}
    \item A volte si osserva un EPSP \textbf{grande} ($\Delta V \approx 1$–$1.2$ mV).
    \item A volte un EPSP \textbf{piccolo} ($\Delta V \approx 0.4$ mV).
    \item A volte si osserva un \textbf{failure} (nessun EPSP: $\Delta V = 0$).
\end{itemize}

La risposta non è deterministica: è \textbf{stocastica}. Ogni stimolazione ha un esito diverso, anche a parità di stimolo presinaptico. Questo è un dato sperimentale robusto e riproducibile.

\subsubsection{Analisi Quantale: Distribuzione Multimodale di $\Delta$V}

La distribuzione degli $\Delta V$ osservati sperimentalmente presenta caratteristicamente più \textbf{picchi (modes) equidistanti}, separati di circa $\mu_1 \approx 0.4$ mV. Questa struttura a picchi multipli è la firma dell'\textbf{analisi quantale} della trasmissione sinaptica.

L'ipotesi fondamentale è:
\begin{itemize}
    \item Il rilascio del neurotrasmettitore avviene in unità discrete dette \textbf{quanti}: ciascun quanto corrisponde alla fusione di \textbf{una singola vescicola} e produce un EPSP di ampiezza $\mu_1$ (EPSP unitario o mEPSP — miniature EPSP).
    \item Il numero di quanti rilasciati per ogni stimolo sinaptico $k$ è una variabile aleatoria.
\end{itemize}



\section{Modello Statistico: Distribuzione Binomiale + Gaussiana}

\subsubsection{La Distribuzione Binomiale per il Numero di Vescicole}

L'ipotesi del modello quantale è che la sinapsi disponga di $N$ siti di rilascio (release sites), ciascuno con una vescicola pronta al rilascio. Per ogni spike presinaptico, ciascun sito rilascia la sua vescicola con probabilità $p$ (probabilità di rilascio), indipendentemente dagli altri siti. Il numero totale di quanti (vescicole) rilasciati $k$ segue quindi una \textbf{distribuzione binomiale}:

\begin{equation}
P(k) = \binom{N}{k} p^k (1-p)^{N-k}
\end{equation}

I casi speciali:
\begin{itemize}
    \item $k = 0$ (failure): $P(0) = (1-p)^N$.
    \item $k = N$ (rilascio massimale): $P(N) = p^N$.
    \item Valore atteso: $\langle k \rangle = Np$ (\textbf{quantal content}).
\end{itemize}

\subsubsection{La Distribuzione Gaussiana per il $\Delta$V dato k}

Fissato il numero di quanti rilasciati $k$, l'ampiezza dell'EPSP $\Delta V$ non è esattamente $k \mu_1$, ma fluttua attorno a questo valore. Questo è dovuto a:
\begin{itemize}
    \item Variabilità nell'ampiezza del singolo EPSP unitario (fluttuazioni del numero di molecole per vescicola, della conduttanza dei recettori, ecc.).
    \item Rumore di fondo nel potenziale di membrana postsinaptico.
\end{itemize}

Si assume che, dato $k$, la distribuzione di $\Delta V$ sia \textbf{gaussiana}:

\begin{equation}
P(\Delta V \mid k) = \mathcal{N}(k\mu_1, \, \sqrt{k}\,\sigma_1)
\end{equation}

dove:
\begin{itemize}
    \item $\mu_1$ è l'ampiezza media del singolo quanto (mEPSP unitario) $\approx 0.4$ mV,
    \item $\sigma_1$ è la deviazione standard dell'EPSP per un singolo quanto.
\end{itemize}

La varianza del $\Delta V$ per $k$ quanti è $k\sigma_1^2$ (perché la somma di $k$ variabili indipendenti con varianza $\sigma_1^2$ ha varianza $k\sigma_1^2$).

\subsubsection{Distribuzione Complessiva del $\Delta$EPSP}

La distribuzione totale degli EPSP osservata sperimentalmente è la \textbf{somma pesata} delle distribuzioni gaussiane (una per ogni $k$), pesata dalla distribuzione binomiale:

\begin{equation}
P(\Delta V) = \sum_{k=0}^{N} P(k) \cdot \mathcal{N}(k\mu_1, \sqrt{k}\,\sigma_1)
\end{equation}

Questa è una \textbf{miscela di gaussiane} (Gaussian mixture model) con pesi dati dalla distribuzione binomiale. Il primo termine ($k=0$) è un delta di Dirac in $\Delta V = 0$ (i failure).

Il fit di questa distribuzione teorica ai dati sperimentali permette di stimare:
\begin{itemize}
    \item $N$: numero di siti di rilascio.
    \item $p$: probabilità di rilascio per sito.
    \item $\mu_1$: ampiezza del quanto unitario.
    \item $\sigma_1$: variabilità del quanto unitario.
\end{itemize}

\begin{tcolorbox}[colback=gray!5, colframe=gray!40, arc=2mm, boxrule=0.5pt]
Questa distribuzione teorica fitta \textbf{perfettamente} i dati sperimentali. Questo costituisce una delle prove più eleganti della natura quantale della trasmissione sinaptica, introdotta originalmente da \textbf{Bernard Katz} (Nobel 1970).
\end{tcolorbox}

\subsubsection{La Trasmissione Sinaptica come Fonte di Stocasticità}

La trasmissione sinaptica è una delle principali \textbf{fonti di stocasticità} nel sistema nervoso. Con $10^{10}$ neuroni e $\sim 10^4$ sinapsi per neurone, il cervello ha $\sim 10^{14}$ sinapsi, ciascuna con una propria probabilità di rilascio. Eppure questo sistema massiccio stocastico produce comportamenti affidabili e riproducibili, grazie a meccanismi di \textbf{averaging di campo medio} (mean-field averaging): le fluttuazioni di una singola sinapsi sono irrilevanti quando centinaia di sinapsi convergono sullo stesso neurone.


\section{Modello Matematico della Conduttanza Sinaptica}

\subsubsection{La Corrente Sinaptica}

La corrente sinaptica prodotta dall'apertura dei canali postsinaptici è descritta dalla forma ohmica:

\begin{equation}
I_{syn}(t) = G_{syn}(t) \cdot (V - E_{syn})
\end{equation}

dove:
\begin{itemize}
    \item $G_{syn}(t) = G_{max} \cdot R(t)$ è la conduttanza sinaptica variabile nel tempo,
    \item $G_{max}$ è la conduttanza massima (quando tutti i canali sono aperti),
    \item $R(t) \in [0, 1]$ è la \textbf{frazione di canali postsinaptici aperti} in funzione del tempo,
    \item $E_{syn}$ è il potenziale di inversione della corrente sinaptica ($E_{syn} = E_{Na} \approx +50$ mV per sinapsi eccitatorie AMPA; $E_{syn} = E_{Cl} \approx -70$ mV per sinapsi inibitorie GABA$_A$).
\end{itemize}

\subsubsection{Equazione Cinetica per R: Un Processo di Markov}

La variabile $R(t)$ descrive la cinetica di transizione dei canali ionici postsinaptici tra lo stato chiuso (C) e lo stato aperto (O):

\begin{equation}
\text{C} \underset{\beta}{\overset{\alpha N(t)}{\rightleftharpoons}} \text{O}
\end{equation}

La dinamica di $R$ come sistema del primo ordine è:

\begin{equation}
\dot{R} = \alpha N(t)(1-R) - \beta R
\end{equation}

dove:
\begin{itemize}
    \item $N(t)$ è la \textbf{concentrazione di neurotrasmettitore} nella fessura sinaptica in funzione del tempo,
    \item $\alpha$ è il rate di transizione chiuso $\to$ aperto (processo bimolecolare),
    \item $\beta$ è il rate di transizione aperto $\to$ chiuso (decadimento monomolecolare),
    \item $(1-R)$ è la frazione di canali nello stato chiuso.
\end{itemize}



Il profilo temporale della concentrazione di neurotrasmettitore $N(t)$ nella fessura sinaptica in risposta a uno spike presinaptico al tempo $t = 0$ è complesso, ma può essere modellato con una \textbf{funzione alpha}:

\begin{equation}
N(t) \propto N_{max} \frac{t}{\tau^2} \cdot e^{-t/\tau_T}
\end{equation}

Questa funzione ha un'ascesa lineare seguita da un decadimento esponenziale, con un picco al tempo $t_{peak} = \tau_T$.


\subsubsection{L'Approssimazione a Scalino (Heaviside)}

Per semplificare la soluzione dell'equazione per $R$, si approssima il profilo temporale del neurotrasmettitore con una \textbf{funzione a scalino di durata finita} $\Delta$:

\begin{equation}
N(t) \approx N_{max} \cdot [\Theta(t)  \Theta(t - \Delta)]
\end{equation}

dove $N_{max}$ è la concentrazione massima di neurotrasmettitore. Questa approssimazione trasforma l'equazione differenziale per $R$ in un sistema \textbf{a coefficienti costanti a tratti}, risolvibile analiticamente in ciascuna delle due fasi.

\subsubsection{Fase di Rising (0 $\leq$ t $\leq$ $\Delta$)}

Durante la fase di rising, $N = N_{max} = \text{cost}$. L'equazione diventa:

\begin{equation}
\dot{R} = (\alpha N_{max} + \beta)\left(\frac{\alpha N_{max}}{\alpha N_{max} + \beta} - R\right)
\end{equation}

La soluzione è un'esponenziale che si avvicina all'asintoto:

\begin{equation}
R_\infty^{rise} = \frac{\alpha N_{max}}{\alpha N_{max} + \beta} \approx 1 \quad \text{(poiché } \alpha N_{max} \gg \beta \text{)}
\end{equation}

con costante di tempo:

\begin{equation}
\tau_{rise} = \frac{1}{\alpha N_{max} + \beta} \approx \frac{1}{\alpha N_{max}}
\end{equation}

Tipicamente $\tau_{rise} \sim 1$ ms, dell'ordine della durata dello spike presinaptico.

\subsubsection{Fase di Decaying (t > $\Delta$)}

Quando $N = 0$ (il neurotrasmettitore è diffuso fuori dalla fessura), l'equazione si riduce a:

\begin{equation}
\dot{R} = -\beta R
\end{equation}

La soluzione è il puro decadimento esponenziale:

\begin{equation}
R(t) = R(\Delta) \cdot e^{-\beta(t - \Delta)}
\end{equation}

con asintoto $R_\infty^{decay} = 0$ e costante di tempo:

\begin{equation}
\tau_{decay} = \frac{1}{\beta}
\end{equation}

La conduttanza sinaptica decade quindi con la costante di tempo $\tau_{decay}$, che è determinata unicamente dal rate di chiusura spontanea $\beta$ dei canali.

\subsubsection{Rappresentazione Continua con la Funzione Alpha}

Invece della soluzione a tratti, si può usare la soluzione dell'equazione completa con la funzione alpha come forzante. Il risultato è una curva con:
\begin{itemize}
    \item Slope iniziale $\sim 1/\tau_{rise}$ durante la fase ascendente,
    \item Decadimento esponenziale con $\tau_{decay}$ dopo il picco.
\end{itemize}

Questa curva a "campana" asimmetrica è la \textbf{forma d'onda tipica} di un EPSP o di una conduttanza sinaptica.

\begin{tcolorbox}[colback=gray!5, colframe=gray!40, arc=2mm, boxrule=0.5pt]
Nel passaggio tra l'approssimazione a scalino e la funzione alpha, il parametro $\Delta$ viene aggiustato in modo che l'\textbf{integrale della corrente} (cioè la carica totale trasferita) sia conservato. Questo garantisce che le due rappresentazioni siano equivalenti dal punto di vista della fisiologia postsinaptica.
\end{tcolorbox}




\subsubsection{La Carica per Singolo Evento Sinaptico}

La carica totale $Q$ consegnata al neurone postsinaptico durante un singolo evento sinaptico è:

\begin{equation}
Q = \int_0^\infty I_{syn}(t) \, dt = G_{max} \int_0^\infty R(t) \cdot (V - E_{syn}) \, dt
\end{equation}

Poiché l'ampiezza del singolo EPSP è piccola ($\Delta V \approx 0.4$ mV $\ll$ soglia $\approx 20$ mV), il potenziale di membrana postsinaptico $V$ rimane quasi costante durante l'evento sinaptico. Si può quindi estrarre $(V - E_{syn})$ dall'integrale:

\begin{equation}
Q \approx G_{max} (V - E_{syn}) \int_0^\infty R(t) \, dt
\end{equation}

Di conseguenza, \textbf{la conservazione della carica equivale alla conservazione dell'integrale di} $R(t)$. Nel passaggio tra diverse approssimazioni del profilo temporale di $N(t)$, si sceglie $\Delta$ (o l'ampiezza della funzione alpha) in modo che $\int R(t) dt$ rimanga invariato.

\subsubsection{Classificazione delle Sinapsi per la Velocità}

Le sinapsi chimiche si classificano in base alle costanti di tempo $\tau_{rise}$ e $\tau_{decay}$:

\begin{table}[htbp]
    \centering
    \renewcommand{\arraystretch}{1.3}
    \begin{tabularx}{\textwidth}{|l|l|l|X|X|}
    \hline
    \textbf{Tipo} & \textbf{$\tau_{rise}$} & \textbf{$\tau_{decay}$} & \textbf{Esempio biologico} & \textbf{Funzione} \\
    \hline
    \textbf{Sinapsi veloci} (fast) & $\sim 1$ ms & $3$–$5$ ms & AMPA, GABA$_A$ & Trasmissione fedele di strutture temporali fini \\
    \hline
    \textbf{Sinapsi lente} (slow) & $\sim 1$–$5$ ms & $10$–$100$+ ms & NMDA, GABA$_B$ & Integrazione della frequenza media del treno \\
    \hline
    \end{tabularx}
\end{table}

Per le sinapsi veloci: $\tau_{rise} \approx \tau_{decay}/5$ o meno, e $\tau_{decay} \sim$ durata dello spike.
Per le sinapsi lente: $\tau_{decay} \gg \tau_{rise}$, e la decaying phase domina il profilo temporale.



\section{Il Filtro Passa-Basso Sinaptico: Sinapsi Veloci e Lente}

\subsubsection{L'Approssimazione Esponenziale per Sinapsi con $\tau_{rise} \ll \tau_{decay}$}

Quando $\tau_{rise} \ll \tau_{decay}$, la fase di rising è trascurabile rispetto alla fase di decaying. In questo caso, si può approssimare la risposta sinaptica con un singolo esponenziale con un piccolo ritardo temporale $\delta_t$ (che incorpora l'effetto della fase ascendente):

\begin{equation}
R(t) \approx A \cdot e^{-(t )/\tau_{decay}} \cdot \Theta(t)
\end{equation}

dove $t_f$ è il tempo dello spike presinaptico.

Questa è l'equazione di un \textbf{filtro del primo ordine} applicato al treno di spike presinaptico. In termini di risposta all'impulso (risposta impulsiva del sistema), la funzione $R(t)$ è un esponenziale decrescente con costante di tempo $\tau_{decay}$.

\subsubsection{Il Filtro Sinaptico come Filtro Passa-Basso}

In termini di elaborazione del segnale, la sinapsi chimica (con $\tau_{rise} \ll \tau_{decay}$) si comporta come un \textbf{filtro passa-basso del primo ordine} con frequenza di taglio:

\begin{equation}
f_c = \frac{1}{2\pi \tau_{decay}}
\end{equation}

Questo significa che:
\begin{itemize}
    \item Componenti del segnale a frequenza $f \ll f_c$ (variazioni lente): passano attraverso il filtro quasi inalterate.
    \item Componenti del segnale a frequenza $f \gg f_c$ (variazioni rapide): vengono attenuate (filtrate).
\end{itemize}

\subsubsection{Sinapsi Veloci: Trasmissione Fedele delle Strutture Temporali}

Per sinapsi con $\tau_{decay} \sim 4$ ms, la frequenza di taglio è:

\begin{equation}
f_c \approx \frac{1}{2\pi \times 4 \times 10^{-3}} \approx 40 \, \text{Hz}
\end{equation}

Questo significa che variazioni del treno di spike a frequenze fino a $\sim 40$ Hz vengono trasmesse fedelmente alla membrana postsinaptica. Poiché le frequenze di firing tipiche nella corteccia sono nell'ordine di decine di Hz, le \textbf{sinapsi veloci riproducono fedelmente le strutture temporali fini} del treno di spike presinaptico.

Il potenziale postsinaptico $V_m$ varia quindi in modo simile al profilo temporale del treno di spike: ogni spike presinaptico produce un EPSP riconoscibile e separato.

\subsubsection{Sinapsi Lente: Integrazione della Frequenza Media}

Per sinapsi con $\tau_{decay} \sim 64$ ms, la frequenza di taglio è:

\begin{equation}
f_c \approx \frac{1}{2\pi \times 64 \times 10^{-3}} \approx 2.5 \, \text{Hz}
\end{equation}

In questo caso, le \textbf{alte frequenze del segnale vengono filtrate}: la corrente postsinaptica non segue i singoli spike, ma è quasi costante e riflette la \textbf{frequenza media} del treno presinaptico. Il potenziale $V_m$ varia attorno a un valore quasi costante con oscillazioni molto smorzate.

Le \textbf{sinapsi lente integrano la frequenza media} del treno di spike: sono sensori di rate code, non di spike code temporale.

\subsubsection{Significato Computazionale}

La coesistenza sull'albero dendritico di \textbf{sinapsi veloci} (AMPA, GABA$_A$) e \textbf{sinapsi lente} (NMDA, GABA$_B$) ha un profondo significato computazionale:

\begin{itemize}
    \item Le sinapsi veloci trasmettono informazione con granularità temporale fine (coding temporale, sincronizzazione tra neuroni, binding).
    \item Le sinapsi lente calcolano la frequenza media (rate coding), permettendo al neurone di integrare l'attività del neurone presinaptico su finestre temporali più lunghe.
\end{itemize}

Il neurone postsinaptico riceve quindi informazioni su \textbf{molteplici scale temporali} simultaneamente, aumentando la sua capacità computazionale.

\begin{tcolorbox}[colback=gray!5, colframe=gray!40, arc=2mm, boxrule=0.5pt]
 Questo è uno dei motivi biologici per cui sotto l'albero dendritico di ogni neurone corticale coesistono sinapsi veloci e lente. Non è ridondanza, è complessità computazionale distribuita.
\end{tcolorbox}

\subsubsection{Condizione per la Perdita di Informazione}

La condizione limite per la perdita completa dell'informazione temporale è:

\begin{equation}
f_c = \frac{1}{\tau_{decay}} < f_{segnale}
\end{equation}

dove $f_{segnale}$ è la frequenza caratteristica del contenuto informativo del treno di spike presinaptico. Se questa condizione è soddisfatta, la sinapsi filtra completamente le variazioni temporali del segnale, e la corrente postsinaptica diventa proporzionale alla sola frequenza media del firing presinaptico.

